\section{The frame: Bessel--Poincar\'e series}\label{sec:frame}

\Cref{st:F} needs the period maps $B_u$ to form a nearly tight frame. The clause that supplies
this is \lvS{} (\lean{InputSG}): for matrices $B_i$ with a common target of dimension
$\abs\alpha$, weights $\gamma_i$, and the normalised Gram blocks
$\Gamma_{ij}=(\gamma_i\gamma_j)^{1/2}B_i^*B_j$,
\begin{align*}
\norm{\Gamma_{ij}}_{\mathrm{HS}}^2&\le(1+\eta_0)\,\frac{d_i\,d_j}{D^*}\quad(i\ne j),&
\norm{\Gamma_{ii}-1}_{\mathrm{HS}}^2&\le e_0\,d_i,&
\abs\alpha&\le(1+\eta_1)\,D^*,
\end{align*}
where $d_i$ is the rank of the $i$-th source. Each map is nearly an isometry, and two different
maps overlap no more than two independent random subspaces of dimensions $d_i$, $d_j$ in a
space of dimension $D^*$ would. This is a nearly \emph{tight} frame and not a nearly
orthonormal family: tightness comes from there being many indices, with $\sum_id_i$ far larger
than $D^*$. Passing from these three bounds to the hypotheses of \cref{st:F} is linear algebra.

\subsection{The maps}
The forms live on a definite unitary group $\mathbf H=\U(W_0)$ of rank $k$ over the host field;
in the application $k$ is the rank parameter of \cref{sec:balance}. The
common target is a space of automorphic forms on an \emph{indefinite} unitary group
$\mathbf G=\U(V)$, of signature $(k+r,\,r+1)$ at each infinite place, with a fixed discrete-series type at
infinity. The two are related by a Bessel subgroup $\mathbf N\rtimes\mathbf H\subset\mathbf G$, attached to an
isotropic flag and a vector of norm $-T(u)$: this is the setting of the Gan--Gross--Prasad
conjectures \cite{ggp2012}. The map $\mathcal P_u$ sends a form on $\mathbf H$ of archimedean type $\pi_u$ to
the Poincar\'e series on $\mathbf G$ built from the Bessel function of the discrete series, and up to
unitary changes of basis and a non-zero scalar the period map $B_u$ is a compression of $\mathcal P_u$.

\subsection{A relative trace formula}
The pairing $\langle\mathcal P_u\phi,\mathcal P_{u'}\phi'\rangle$ is expanded as a convergent sum of operators
indexed by integral double cosets, each an archimedean orbital integral \lvS{}
(\lean{KExpansion}); the terms that are not \emph{relevant} vanish. The three bounds are read
off this sum.
\begin{itemize}
\item \emph{The identity term.} For a term whose representative lies in the unipotent group
  $\mathbf N$, the operator is a scalar \lvS{} (\lean{oot_eq}): a character value times
  $\mathrm{vol}(\mathbf N\text{-domain})\cdot\mathrm{vol}(\mathbf H)/\dim\pi_u$. The normalisation of the maps
  makes this $1$.
\item \emph{Different indices.} If the archimedean parameters of $u$ and $u'$ differ at every
  place, then for a relevant degenerate term every intertwining operator between the two
  representations vanishes \lvS{} (\lean{SchurNFC}); so such terms contribute nothing off the
  diagonal.
\item \emph{Uniqueness of Bessel models.} At the archimedean place, in a range where the
  parameters are separated by a gap of at least $120\,(2r+k+1)^3$, any two Bessel maps from a
  discrete series of $\U(k+r,r+1)$ to $\pi_u$ are proportional on smooth vectors \lvS{}
  (\lean{UniqGap}); compare \cite{sunzhu2012,jsz2010}.
\item \emph{The Siegel--Weil formula} is used in its convergent range, for frames of fewer
  vectors than half the dimension of the Hermitian space, and as a mass
  formula \lvS{} (\lean{SiegelDensities}): the mean over $[\U(V)]$ of the theta series of
  frames with a given non-degenerate Gram matrix is the product of the local densities. The
  main term of the off-diagonal blocks is derived from it and from the uniqueness of Bessel
  models (\lean{rowMainEqP_of_siegelWeil}).
\end{itemize}
The outcome for two different indices is \lvS{} (\lean{pbApplied}): for $k$ large, if the
rows of the expansion that are large or not scalar contribute at most
$\exp(-c\,kf\log k)$ times the main term, and the non-generic rows at most
$2^{-(k-2r-5)}\exp\bigl(C(2r+2)^5f\log b\bigr)$ times it, then
\[
\norm{\Gamma_{uu'}}^2\ \le\ \bigl(1+b^{-6}\bigr)\,\frac{d_u\,d_{u'}}{D^*}.
\]
Here $f=\abs\Phi$ is the number of infinite places of the host. The decay is driven by the rank $k$.

Every ingredient is standard: unfolding as in the Petersson formula, Schur orthogonality,
decay of matrix coefficients, uniqueness of Bessel models, Siegel--Weil as a mass formula,
lattice-point counting. What appears to be new is the uniformity: the bound holds for the whole
family of archimedean types at once, with every off-diagonal block within a factor $1+b^{-6}$
of the random-subspace value.

The abstract maps to which \cref{st:F} is applied are these period maps, up to linear
isometries and a non-zero scalar \lvS{} (\lean{BesselId}).
